nommen haben; bis dahin hätte aber die Tafel über 21000 Glieder. Was aber noch wichtiger ist: man darf keinesweges die dritten Differenzen überall als verschwindend betrachten, wie \textsc{Leonelli} sich einbildet, der nur auf die zweiten Rücksicht zu nehmen für hinreichend hält. Zu diesem Irrthume scheint ihn die von ihm berechnete Probe des Anfangs der Tafel verleitet zu haben, wo freilich die dritten Differenzen verschwinden; allein eine leichte Rechnung zeigt, dass diese dritten Differenzen weiterhin \emph{zunehmen} und allerdings bedeutend werden können, wenn man 14 Decimalen geben, und keine kleinern Differenzen bei $P$, als 0,0005 gebrauchen will: für $x = 2+\BareRoot{3}$, oder für $P = 0,5719$, wo die dritten Differenzen ihr Maximum erreichen, finden wir unter obigen Voraussetzungen ihren Werth $= 6377$ Einheiten in der vierzehnten Decimale. Weit entfernt also, dass \textsc{Leonelli} mit weniger als 5000 Gliedern ausreichen könnte, müsste die Tafel, wenn die zweiten Differenzen überall hinlänglich sein sollen, eine so grosse Ausdehnung erhalten, dass ihre Berechnung die Mühe keinesweges belohnen würde. Allein wozu auch vierzehn Decimalen? Rechnungen, wo eine solche Schärfe nöthig wäre, kommen ja nur höchst selten vor, und für einen so seltenen Fall eine doch nicht sehr bedeutende Abkürzung der Arbeit zu erhalten, daran ist wenig gelegen. Hingegen für Rechnungen, die täglich vorkommen, und wo sieben Decimalen völlig hinreichen, würde eine, wenn auch nur mässige, aber oft wiederkommende Erleichterung der Arbeit allerdings schätzbar sein. Dann müsste aber die Tafel, um auch die nöthige Bequemlichkeit zu gewähren, so eingerichtet werden, dass überall auch die zweiten Differenzen verschwinden; mit etwas mehr als 15000 Gliedern liesse sich diess bequem erreichen, die nur ein mässiges Bändchen machen würden. Doch zu einer weitern Ausführung ist hier nicht der Ort.

Noch eine Probe, wie oberflächlich \textsc{Leonelli} seinen Gegenstand behandelt hat, gibt die S.\,60 vorgetragene Formel, für einen vorgegebenen Werth von $R$, der zwischen zwei der Tafel fällt, den entsprechenden von $Q$ mit Rücksicht auf die zweiten Differenzen zu finden. Diese Formel, deren Deduction dem Übersetzer, seinem Geständnisse zufolge, so viele vergebliche Mühe gemacht hat, ist ganz falsch; statt
\[
2\delta(r'n - n) \quad \text{sollte nemlich stehen} \quad \frac{(D+d)\delta rn}{d}.
\]
